% Progress report for the supervisor, 2026-10-17.
% Build from this directory:  tectonic 2026-10-17-progress-report.tex
% Figures are the PNGs in graphs/ at the repo root (exported from notebook 14).
\documentclass[10pt,a4paper]{article}

\usepackage[a4paper,margin=2cm]{geometry}
\usepackage{graphicx}
\usepackage{xcolor}
\usepackage{caption}
\usepackage{enumitem}
\usepackage{titlesec}
\usepackage{float}
\usepackage{listings}
\usepackage{booktabs}
\usepackage{tabularx}
\usepackage[hidelinks]{hyperref}

\graphicspath{{../../graphs/}}

\definecolor{ink2}{HTML}{52514E}
\captionsetup{font={small,color=ink2},labelfont={small,color=ink2},justification=raggedright,singlelinecheck=false}
\titleformat{\section}{\large\bfseries}{\thesection.}{0.5em}{}
\titlespacing*{\section}{0pt}{1.2em}{0.4em}
\setlist{itemsep=0.25em,topsep=0.3em}
\setlength{\parindent}{0pt}
\setlength{\parskip}{0.5em}
\setlength{\intextsep}{8pt}
\captionsetup{skip=3pt}
\renewcommand{\arraystretch}{1.25}
\newcolumntype{L}{>{\raggedright\arraybackslash}X}

\definecolor{codebg}{HTML}{F6F6F4}
\lstset{basicstyle=\ttfamily\footnotesize,breaklines=true,columns=fullflexible,keepspaces=true,
        backgroundcolor=\color{codebg},frame=single,rulecolor=\color{ink2!40},
        xleftmargin=4pt,xrightmargin=4pt,aboveskip=6pt,belowskip=6pt}

\hypersetup{pdftitle={Supplement KG progress report},pdfauthor={Yanfeng Li}}

% One figure from graphs/; the optional first argument is its width (default: full width).
\newcommand{\fig}[3][\linewidth]{%
  \begin{figure}[H]
    \centering
    \includegraphics[width=#1]{#2}
    \caption{#3}
  \end{figure}}

\begin{document}

{\LARGE\bfseries Supplement knowledge graph: progress report}\\[0.3em]
{\small\color{ink2} Yanfeng Li \textperiodcentered{} 17 October 2026}

\section*{Summary}
\begin{itemize}
  \item Purpose and risk claims have been extracted from \textbf{all 214,780 DSLD labels}, with no
        failures: 466,630 purpose items and 355,134 risk items, about 98\% copied verbatim from the label.
  \item The main run \textbf{finished on 15 October} after \textbf{42.6 hours} (13 Oct 18:22 to
        15 Oct 13:00 UTC), averaging \textbf{74 generated tokens/s} overall with 4 requests in
        parallel (18.6 tokens/s per request). It got 214,773 labels right first time. The other 7 had
        outputs too long for the output-token limit, so on 16 October they were rerun with a longer
        limit, and all 7 came back ok.
  \item Next, the extraction will be evaluated for faithfulness and correctness, with two LLM judges
        and a human reviewer.
\end{itemize}

\section{Extraction results}

A local LLM (Qwen, run on our own GPU) read each DSLD label and extracted two kinds of statements:
\textbf{purposes} (a claimed effect or use) and \textbf{risks} (warnings, contraindications,
interactions, side effects). Each item is meant to be a verbatim piece of the label, so that it can be
checked against the source and used on its own by a later NER step. The full prompt is in
Appendix~\ref{app:prompt}.

\fig{02_extraction_groups.png}{Labels grouped by what was extracted. The labels with neither are
mostly very short (median 26 words).\label{fig:groups}}

\fig{03_faithfulness.png}{Faithfulness to the label, checked automatically on every item.\label{fig:faith}}

The label text is split into \textbf{segments}, one per statement or ingredient note. Each item falls
into one of four types:
\begin{itemize}
  \item \textbf{Verbatim:} the item appears in the label exactly, character for character.
  \item \textbf{Stitched within a segment:} not verbatim, but every word comes from the same
        segment. The model built it from words of one label sentence, leaving out some words in between.
  \item \textbf{Spread over segments:} every word is in the label, but taken from more than one segment.
  \item \textbf{Invented words:} the item contains at least one word that is not in the label.
\end{itemize}

Stitched items usually come from the model taking one benefit out of a list, as the prompt asks it to
put one claim per line:

\begin{table}[H]
\small
\begin{tabularx}{\linewidth}{@{}>{\raggedright\arraybackslash}p{0.33\linewidth} L@{}}
\toprule
Extracted item & Label segment it came from (left-out words in bold) \\
\midrule
promotes heart and circulatory health & promotes \textbf{sugar metabolism} heart and circulatory health \\
support vision & \ldots{} reduce oxidative stress and support \textbf{healthy cell development
  cardiovascular function and} vision \ldots \\
increase energy & \ldots{} may help you increase \textbf{the body s metabolic rate and} energy \ldots \\
\bottomrule
\end{tabularx}
\end{table}

These keep the label's meaning. Occasionally stitching changes it: one item reads ``feel the euphoria
turns into a feeling of intensity aggression and rage'', joining words from far apart in a long
marketing paragraph. So a stitched item is not necessarily wrong, but it cannot be verified
automatically. Whether its meaning was kept is checked in the evaluation (Section~\ref{sec:eval}).

Being verbatim shows that the model almost never invents text. It does not show that each item is a
correct claim in the right bucket. That is what the evaluation measures.

\section{Evaluation plan}\label{sec:eval}

The evaluation checks two things: whether each item is \textbf{faithful} to the label, and whether it
is in the \textbf{right bucket} (a real purpose, a real risk, or neither).

\begin{itemize}
  \item \textbf{Precision} (per item): from a sample stratified by the four item types above
        (verbatim, stitched, spread over segments, invented words), 100 items per type for purposes and
        for risks, reweighted afterwards (Figure~\ref{fig:itemsample}). A plain random sample would be
        almost all verbatim, so the rarer types could not be measured. Errors are split into \emph{not a claim}, \emph{wrong
        bucket} and \emph{unfaithful}.
  \item \textbf{False negative and false positive rates} (per label): whether each label's purposes and
        risks were missed, or extracted when the label has none. From a sample stratified by the four
        groups in Figure~\ref{fig:groups}, 150 labels each, reweighted to the catalog afterwards. A label's group
        decides which errors it can show (Figure~\ref{fig:labelsample}).
  \item Every rate is reported with a 95\% confidence interval.
\end{itemize}

\fig{08_eval_item_sampling.png}{Stratified item sample. Equal samples per type make sure the rare
types, which most need checking, are well covered.\label{fig:itemsample}}

\fig{06_eval_sampling.png}{Stratified label sample. Equal samples per group make sure the small
groups, where missed claims can show up, are well covered.\label{fig:labelsample}}

\fig{07_eval_judging_flow.png}{Judging workflow.}

Two LLM judges label each item independently. I decide every case where they disagree, and I audit a
random 10\% of the cases where they agree, labelling first without seeing their verdicts. The audit
tells how far their agreement can be trusted. A judge that says a claim was missed has to quote the
label, and the quote is checked automatically. A 50-item pilot fixes the rubric first. Judge
reliability is reported as Cohen's $\kappa$.

\appendix
\section{Extraction prompt}\label{app:prompt}

The system prompt used for the full run, verbatim. The label text is given to the model as the
user message, and the model answers in the line format at the end.

\lstinputlisting[basicstyle=\ttfamily\fontsize{7.5}{9}\selectfont]{../../engine/repositories/prompts/purpose_risk/prompt.md}

\end{document}
